% ==========================================================================
%  Chapter 88 -- Validation Reboot Master Plan
% ==========================================================================

\chapter{Validation Reboot Master Plan}
\label{ch:validation-reboot-master-plan}

\paragraph{Normative status.}
This chapter defines the validation program that shall replace the current
legacy mixture of exploratory drivers, partial regression surfaces, and
large executables that do not yet provide a scientifically clean entry point.
Its purpose is methodological before it is numerical: to prevent the library
from inheriting validation authority from surfaces that were never designed as
controlled benchmarks.

\section{Why the validation campaign must restart from zero}

The library already contains many tests and several large validation-oriented
drivers.  That is valuable, but not sufficient.  A test only proves what it
actually constrains, and a driver only validates what its inputs, outputs, and
acceptance criteria make explicit.  If those criteria are weak, hidden, or
historically accidental, the resulting evidence cannot be promoted to a
scientific claim without re-audit.

The new validation program therefore adopts the following rule:
\begin{equation}
  \text{claim}
  \;\Longrightarrow\;
  \text{computational artifact}
  \;\Longrightarrow\;
  \text{audited evidence}
  \;\Longrightarrow\;
  \text{acceptance gate}.
\end{equation}
The campaign is not allowed to skip any arrow in this chain.

This rule follows directly from the continuum and finite-element structure
already fixed elsewhere in the project.  The mechanical object is defined in
terms of configuration, work-conjugate pairs, and history variables; the
discrete object is defined in terms of residual, tangent, and incremental state
commitment; and the validation object must therefore be defined in terms of a
traceable relation between those theoretical statements and reproducible
numerical evidence.  See, in particular,
\cite{BonetWood2008,Crisfield1991_NLFEA_Vol1,DeBorst2012_NonlinearFEM}.

\section{Legacy quarantine policy}

The restart does \emph{not} imply immediate destructive deletion.  Current
surfaces remain useful as audited input, but they are no longer allowed to act
as normative validation anchors until the replacement baseline exists.

\begin{table}[htbp]
  \centering
  \caption{Legacy validation surfaces during the reboot.}
  \label{tab:validation-legacy-policy}
  \small
  \begin{tabular}{@{}p{0.30\textwidth}p{0.23\textwidth}p{0.37\textwidth}@{}}
    \toprule
    \textbf{Legacy surface} & \textbf{Temporary status} & \textbf{Disposition rule} \\
    \midrule
    \texttt{src/validation/ComprehensiveCyclicValidation.cpp}
    & audited input only
    & keep available for provenance; quarantine to \texttt{.old} only after the
      new single-column baseline is operational. \\

    \texttt{src/validation/TableCyclicValidationStructural.cpp}
    & audited input only
    & useful as seed data because it already traverses beam columns from
      \(N=2\) to \(N=10\), but it may not act as validated truth. \\

    \texttt{doc/ch82\_cyclic\_validation.tex} and
    \texttt{doc/ch85\_comprehensive\_cyclic\_multiscale\_validation.tex}
    & historical record
    & preserve as provenance; do not present as the normative validation
      narrative once the rebooted campaign is in place. \\
    \bottomrule
  \end{tabular}
\end{table}

\section{First scientific target}

The first normative target is deliberately narrow:
\begin{quote}
  \emph{a single reinforced-concrete rectangular column subjected to
  progressively amplified cyclic lateral displacement up to approximately
  \(50\) to \(100\) mm, with optional axial compression, first in one lateral
  direction and later, if justified, in two directions.}
\end{quote}

This target is small enough to be inspectable, rich enough to expose section
nonlinearity, hysteresis, stiffness degradation, and eventually crack-driven
continuum effects, and disciplined enough to support direct comparison between
reduced structural models and full continuum models.

\section{Canonical campaign phases}

The new campaign is frozen in code through
\texttt{src/validation/ValidationCampaignCatalog.hh}.  Table
\ref{tab:validation-reboot-phases} summarizes the representative workstreams.

\begin{table}[htbp]
  \centering
  \caption{Canonical validation reboot workstreams.}
  \label{tab:validation-reboot-phases}
  \scriptsize
  \setlength{\tabcolsep}{3pt}
  \begin{tabularx}{\textwidth}{@{}p{0.18\textwidth}p{0.12\textwidth}p{0.14\textwidth}p{0.25\textwidth}X@{}}
    \toprule
    \textbf{Workstream} & \textbf{Phase} & \textbf{Priority} &
    \textbf{Primary objective} & \textbf{Gate impact} \\
    \midrule
    Governance reset and evidence protocol
    & 0 & mandatory
    & Freeze the rule that legacy surfaces may inform the campaign but cannot
      anchor validation without re-audit.
    & Blocks every later gate. \\

    Formulation and solver scope audit
    & 1 & mandatory
    & Re-audit TL, UL, corotational, Newton, dynamics, and continuation as
      computational slices.
    & Blocks structural, continuum, and escalation stages. \\

    Uniaxial material and section baseline
    & 2 & mandatory
    & Re-audit steel, concrete, and fiber-section ingredients before using them
      in the reference column.
    & Blocks both reduced and continuum column campaigns. \\

    Reduced-order RC column matrix
    & 3 & mandatory
    & Validate the rectangular column with Timoshenko beam models from
      \(N=2\) to \(N=10\), including hysteresis and moment-curvature outputs.
    & First structural acceptance gate. \\

    Beam integration family extension
    & 3 & mandatory
    & Replace the current one-rule beam integration path with an explicit
      Gauss/Lobatto/Radau strategy family.
    & Required before the structural column campaign can close honestly. \\

    Continuum RC column matrix
    & 4 & mandatory
    & Build the equivalent 3D continuum column with Hex8/20/27 and consistent
      top-face displacement control.
    & Blocks continuum closure and escalation. \\

    Reduced-versus-continuum equivalence
    & 5 & mandatory
    & Quantify agreement in stiffness, strength, dissipation, and loop shape
      before reopening larger structural examples.
    & Final gate before large structures. \\

    Reinforcement discretization extension
    & 4 & conditional
    & Introduce \texttt{TrussElement<Nnodes>} or an equivalent richer
      reinforcement carrier only if the bonded 2-node path proves insufficient.
    & Does not block the first campaign by default. \\

    Alternative concrete model extension
    & 2 & conditional
    & Add other concrete models only if the audited campaign shows that current
      uniaxial or continuum concrete paths are not sufficient.
    & Evidence-driven, not promise-driven. \\

    Force-based structural element path
    & 6 & deferred
    & Open a later growth track toward flexibility-based/force-based elements.
    & Valuable, but intentionally non-blocking for the first column baseline. \\
    \bottomrule
  \end{tabularx}
\end{table}

\section{Dependency graph of the reboot}

\begin{figure}[htbp]
  \centering
  \begin{tikzpicture}[
      node distance=8mm and 12mm,
      box/.style={draw, rounded corners, align=center, fill=blue!5,
                  minimum width=34mm, minimum height=9mm},
      opt/.style={draw, rounded corners, align=center, fill=orange!10,
                  minimum width=34mm, minimum height=9mm},
      arr/.style={-{Latex[length=2.5mm]}, thick}
    ]
    \node[box] (gov) {Phase 0\\Governance};
    \node[box, right=of gov] (scope) {Phase 1\\Formulation + solver audit};
    \node[box, right=of scope] (mat) {Phase 2\\Materials + sections};

    \node[box, below=of scope] (beam) {Phase 3\\Reduced RC column};
    \node[box, right=of beam] (cont) {Phase 4\\Continuum RC column};
    \node[box, right=of cont] (eqv) {Phase 5\\Cross-model equivalence};
    \node[box, below=of eqv] (full) {Phase 6\\Large-structure escalation};

    \node[opt, below=of beam] (rebar) {Conditional\\TrussElement<N> / richer rebar};
    \node[opt, left=of rebar] (conc) {Conditional\\alternative concrete models};
    \node[opt, below=of cont] (force) {Deferred\\force-based elements};

    \draw[arr] (gov) -- (scope);
    \draw[arr] (scope) -- (mat);
    \draw[arr] (mat) -- (beam);
    \draw[arr] (beam) -- (cont);
    \draw[arr] (cont) -- (eqv);
    \draw[arr] (eqv) -- (full);
    \draw[arr] (beam) -- (rebar);
    \draw[arr] (mat) -- (conc);
    \draw[arr] (eqv) -- (force);
  \end{tikzpicture}
  \caption{Dependency structure of the validation reboot.  The first column
  baseline is the mainline; richer reinforcement carriers, additional concrete
  models, and force-based elements are controlled side tracks.}
  \label{fig:validation-reboot-dependency-graph}
\end{figure}

\section{Beam-axis quadrature family as a structural validation enabler}

One concrete blocker identified during the reboot is subtle but important:
for a section-based beam formulation, the axial quadrature rule is not merely
an integration detail. It also defines the \emph{section stations} at which
constitutive history lives and, for the current mixed Timoshenko formulation,
the tying stations used by the reduced shear field. Changing the beam-axis rule
therefore changes both
\begin{enumerate}
  \item the accuracy of the element integral, and
  \item the discrete placement of sectional history.
\end{enumerate}
This is precisely why the reboot could not honestly validate a rectangular RC
column while the beam path remained nominally tied to one quadrature family.

The library now exposes an explicit compile-time beam-axis quadrature family in
\texttt{src/numerics/numerical\_integration/BeamAxisQuadrature.hh}.  The
current family includes:
\begin{itemize}
  \item Gauss--Legendre rules, exact for polynomials up to degree
    \(2m-1\) with \(m\) stations;
  \item Gauss--Lobatto rules, endpoint-inclusive and exact up to degree
    \(2m-3\);
  \item Gauss--Radau rules, single-endpoint-inclusive and exact up to degree
    \(2m-2\).
\end{itemize}
These exactness properties are classical
\cite{Crisfield1991_NLFEA_Vol1,bathe2014_FiniteElementProcedures} and matter
here because the reduced-column campaign explicitly compares section-placement
families as part of the model, not as a hidden numerical afterthought.

For the degenerate one-station case (\(m=1\)), used by the 2-node beam family,
the new API adopts a midpoint extension for Lobatto and Radau in order to keep
the beam-axis contract uniform without introducing a pathological endpoint-only
rule into the baseline structural path. This is a computational convenience,
not a claim about classical one-point Lobatto/Radau theory, and it is
documented as such.

Architecturally, the key point is that the hot path remains static. The
quadrature family is selected at compile time, and
\texttt{TimoshenkoBeamN<N>} reconstructs its reduced shear interpolation from
the actual station coordinates exposed by the bound geometry rule. No runtime
dispatch is inserted in the quadrature loop.

\begin{lstlisting}[language=C++,caption={Compile-time beam-axis quadrature selection used by the rebooted structural campaign.},label={lst:beam-axis-quadrature-selector}]
using GL3  = BeamAxisQuadratureT<BeamAxisQuadratureFamily::GaussLegendre, 3>;
using GLob = BeamAxisQuadratureT<BeamAxisQuadratureFamily::GaussLobatto, 3>;
using GRad = BeamAxisQuadratureT<BeamAxisQuadratureFamily::GaussRadauLeft, 3>;

static_assert(BeamAxisQuadratureTraits<
    BeamAxisQuadratureFamily::GaussLegendre, 3>::polynomial_exactness_degree == 5);
static_assert(BeamAxisQuadratureTraits<
    BeamAxisQuadratureFamily::GaussLobatto, 3>::includes_left_endpoint);
\end{lstlisting}

This enabling layer is now frozen by
\texttt{tests/test\_beam\_axis\_quadrature.cpp}.  The regression constrains:
\begin{itemize}
  \item monomial exactness of the three rule families;
  \item mirror consistency of left/right Radau;
  \item reuse of the same \texttt{TimoshenkoBeamN} formulation with Lobatto
    and Radau section stations;
  \item and symmetry of the resulting structural tangent.
\end{itemize}
Methodologically, that is the correct precondition for the next structural
validation step: the column campaign may now compare integration families as an
explicit modeling choice rather than inheriting one silently from the element
implementation.

\section{Reduced RC column structural matrix}

The next layer of the reboot is no longer a generic workstream label but a
concrete computational matrix.  For the reduced reinforced-concrete column, the
current structural campaign is organized around the typed slice
\begin{equation}
  \mathcal{S}_{N,q,\mathcal{K}}
  \;=\;
  \texttt{TimoshenkoBeamN<N>}
  \;+\;
  q
  \;+\;
  \mathcal{K},
\end{equation}
where \(N \in \{2,\dots,10\}\) is the beam interpolation order,
\(q \in \mathcal{Q}_{\text{beam}}\) is the beam-axis quadrature family, and
\(\mathcal{K}\) is the beam kinematic/formulation slice under audit.

The quadrature family set is now
\begin{equation}
  \mathcal{Q}_{\text{beam}}
  =
  \left\{
    \text{Gauss--Legendre},
    \text{Gauss--Lobatto},
    \text{Gauss--Radau}_{L},
    \text{Gauss--Radau}_{R}
  \right\},
\end{equation}
while the formulation axis retained in the matrix is
\begin{equation}
  \mathcal{K}
  \in
  \left\{
    \text{small strain},
    \text{corotational},
    \text{Total Lagrangian},
    \text{Updated Lagrangian}
  \right\}.
\end{equation}
The crucial point is that this axis is kept \emph{even when the runtime path
does not exist}.  Otherwise the validation program would silently collapse the
set of desirable claims to the subset that happens to be implemented today.

This structural matrix now lives canonically in
\texttt{src/validation/ReducedRCColumnStructuralMatrixCatalog.hh}.  Its present
reading is intentionally conservative and can be summarized as follows:
\begin{equation}
  \mathfrak{B}_{\text{phase3}}
  =
  \left\{
    \mathcal{S}_{N,q,\text{small strain}}
    \;:\;
    N=2,\dots,10,\;
    q \in \mathcal{Q}_{\text{beam}}
  \right\},
\end{equation}
that is, \(36\) runtime-ready baseline cases.
The two remaining subsets are not hidden:
\begin{align}
  \mathfrak{E}_{\text{CR}}
  &=
  \left\{
    \mathcal{S}_{N,q,\text{corotational}}
  \right\},
  &
  |\mathfrak{E}_{\text{CR}}| &= 36,
  \\
  \mathfrak{U}_{\text{finite}}
  &=
  \left\{
    \mathcal{S}_{N,q,\text{TL}},
    \mathcal{S}_{N,q,\text{UL}}
  \right\},
  &
  |\mathfrak{U}_{\text{finite}}| &= 72.
\end{align}
Here \(\mathfrak{E}_{\text{CR}}\) denotes a genuine \emph{family extension}
frontier: beam corotational kinematics exist elsewhere in the library, but not
yet as a \texttt{TimoshenkoBeamN<N>} family path.  By contrast,
\(\mathfrak{U}_{\text{finite}}\) denotes a \emph{missing formulation family}:
neither TL nor UL is currently a runtime beam formulation for this family, so
claiming them would require more than a local patch.

\begin{table}[htbp]
  \centering
  \caption{Current structural-matrix interpretation for the reduced RC column.}
  \label{tab:reduced-rc-column-structural-matrix}
  \small
  \begin{tabular}{@{}p{0.21\textwidth}p{0.15\textwidth}p{0.14\textwidth}p{0.15\textwidth}p{0.23\textwidth}@{}}
    \toprule
    \textbf{Slice family} & \textbf{Case count} & \textbf{Runtime path} &
    \textbf{Phase-3 baseline} & \textbf{Meaning} \\
    \midrule
    \(\mathcal{S}_{N,q,\text{small strain}}\)
    & 36 & yes & yes
    & Current audited reduced-column baseline. \\

    \(\mathcal{S}_{N,q,\text{corotational}}\)
    & 36 & no & no
    & Planned \texttt{TimoshenkoBeamN} kinematic extension; not yet a runtime
      path. \\

    \(\mathcal{S}_{N,q,\text{TL}}\),
    \(\mathcal{S}_{N,q,\text{UL}}\)
    & 72 & no & no
    & Not available for the present beam family; requires a distinct
      finite-kinematics beam formulation path. \\
    \bottomrule
  \end{tabular}
\end{table}

\begin{figure}[htbp]
  \centering
  \begin{tikzpicture}[
      node distance=8mm and 12mm,
      box/.style={draw, rounded corners, align=center, fill=blue!5,
                  minimum width=34mm, minimum height=9mm},
      ext/.style={draw, rounded corners, align=center, fill=orange!10,
                  minimum width=34mm, minimum height=9mm},
      miss/.style={draw, rounded corners, align=center, fill=red!8,
                   minimum width=34mm, minimum height=9mm},
      arr/.style={-{Latex[length=2.5mm]}, thick}
    ]
    \node[box] (catalog) {ReducedRCColumn\\StructuralMatrixCatalog};
    \node[box, right=of catalog] (baseline) {ReducedRCColumn\\StructuralBaseline};
    \node[box, right=of baseline] (legacy) {Legacy wrapper\\TableCyclicValidationStructural};

    \node[ext, below=of baseline] (cr) {Corotational\\family extension};
    \node[miss, below=of cr] (finite) {TL / UL beam\\family absent};

    \draw[arr] (catalog) -- (baseline);
    \draw[arr] (baseline) -- (legacy);
    \draw[arr] (catalog) -- (cr);
    \draw[arr] (catalog) -- (finite);
  \end{tikzpicture}
  \caption{The new reduced-column layer separates the current runtime baseline
  from two different kinds of missing scope: a planned corotational extension
  and a genuinely absent finite-kinematics beam family.}
  \label{fig:reduced-rc-column-structural-matrix}
\end{figure}

This split is not merely descriptive.  The runtime baseline is now also
available through
\texttt{src/validation/ReducedRCColumnStructuralBaseline.hh/.cpp}, so the
legacy beam-column case no longer owns the only implementation of the reduced
column path.

\begin{lstlisting}[language=C++,caption={The reduced-column structural baseline now returns both global hysteresis and section-level observables through a clean runtime API.},label={lst:reduced-rc-column-structural-baseline-api}]
using namespace fall_n::validation_reboot;

auto result = run_reduced_rc_column_small_strain_beam_case_result(
    ReducedRCColumnStructuralRunSpec{
        .beam_nodes = 3,
        .beam_axis_quadrature_family = BeamAxisQuadratureFamily::GaussLobatto,
        .axial_compression_force_mn = 0.02,
        .write_hysteresis_csv = false,
        .write_section_response_csv = false,
        .print_progress = false,
    },
    out_dir,
    cfg);

auto& hysteresis = result.hysteresis_records;
auto& base_side_section_history = result.section_response_records;
\end{lstlisting}

The corresponding regression,
\texttt{tests/test\_reduced\_rc\_column\_structural\_matrix.cpp}, freezes both
the compile-time counts and a small runtime smoke pass over the new baseline.
That smoke pass is intentionally modest: it does \emph{not} claim physical
validation yet.  It only proves that the new structural baseline can traverse
the same reduced-column path under the three audited quadrature families
without relying on the legacy monolithic driver.

\section{Reduced RC column claim--artifact--evidence matrix}

Once the structural matrix exists, one more distinction becomes necessary.
Knowing that a runtime slice is available does not yet answer what we are
scientifically allowed to claim about that slice.  For the reduced RC column,
the reboot therefore introduces an additional claim layer in
\texttt{src/validation/ReducedRCColumnValidationClaimCatalog.hh}.

Formally, the phase-3 reduced-column claim set can be written as
\begin{equation}
  \mathfrak{C}_{\text{RC},\,\mathrm{phase3}}
  =
  \left\{
    C_q,\,
    C_{\mathrm{rt}},\,
    C_{N},\,
    C_{\sigma_N},\,
    C_{\mathrm{hys}},\,
    C_{M\kappa},\,
    C_{\mathrm{qsens}}
  \right\},
\end{equation}
where:
\begin{align}
  C_q &:\text{ the beam-axis quadrature family is an explicit modelling axis,}
  \\
  C_{\mathrm{rt}} &:\text{ the runtime beam-column surface exists for } N=2,\dots,10,
  \\
  C_N &:\text{ a node-refinement convergence claim can eventually be posed,}
  \\
  C_{\sigma_N} &:\text{ an optional axial-compression path exists,}
  \\
  C_{\mathrm{hys}} &:\text{ a base-shear-vs-drift hysteresis observable exists,}
  \\
  C_{M\kappa} &:\text{ a base-side moment-curvature observable exists at the active}
                 \text{ section station closest to the fixed end,}
  \\
  C_{\mathrm{qsens}} &:\text{ a quadrature-sensitivity claim can eventually be posed.}
\end{align}
The methodological point is that these claims are \emph{not} all at the same
readiness level.

For the reduced beam-column baseline, the moment-curvature observable is now
defined explicitly as
\begin{equation}
  \mathcal{O}_{M\kappa}^{\mathrm{base}}
  :=
  \left\{
    \bigl(\kappa_y^{(g_\star)}(t_n),\,M_y^{(g_\star)}(t_n)\bigr)
  \right\}_{n=0}^{N_t},
  \qquad
  g_\star = \arg\min_g \xi_g,
\end{equation}
where \(\xi_g\) denotes the beam-axis coordinate of section station \(g\).
This definition is deliberate: for Lobatto and left-Radau rules,
\(g_\star\) may coincide with the boundary section, whereas for
Gauss--Legendre and right-Radau it is the nearest active station on the fixed
side.  The observable therefore exists cleanly at runtime, but its physical
interpretation still remains coupled to the quadrature-sensitivity benchmark.

\begin{table}[htbp]
  \centering
  \caption{Claim--artifact--evidence matrix for the reduced RC column baseline
  before physical benchmarking.}
  \label{tab:reduced-rc-column-claim-trace}
  \scriptsize
  \setlength{\tabcolsep}{3pt}
  \begin{tabularx}{\textwidth}{@{}p{0.22\textwidth}p{0.18\textwidth}p{0.18\textwidth}p{0.12\textwidth}X@{}}
    \toprule
    \textbf{Claim} & \textbf{Artifact} & \textbf{Current evidence} &
    \textbf{Readiness} & \textbf{Next gate} \\
    \midrule
    Explicit beam-axis quadrature family
    & \texttt{BeamAxisQuadrature.hh}
    & exactness and station-symmetry regression
    & frozen enabler
    & none; this is an enabling contract, not a physical benchmark \\[0.4mm]

    Runtime reduced-column surface for \(N=2,\dots,10\)
    & \texttt{ReducedRCColumnStructuralBaseline}
    & matrix counts + runtime smoke
    & runtime-ready
    & node-refinement benchmark suite \\[0.4mm]

    Optional axial-compression path
    & \texttt{ReducedRCColumnStructuralBaseline}
    & finite axial-preload smoke run
    & runtime-ready
    & axial-load sensitivity suite \\[0.4mm]

    Hysteresis observable
    & \texttt{hysteresis.csv} writer contract
    & output-contract regression
    & runtime-ready
    & hysteresis-shape / dissipation benchmark suite \\[0.4mm]

    Base-side moment-curvature observable
    & \texttt{ReducedRCColumnStructuralBaseline} +
      \texttt{moment\_curvature\_base.csv}
    & output-contract regression for the active base-side section station
    & runtime-ready
    & section benchmark suite + quadrature-sensitivity interpretation \\[0.4mm]

    Node-refinement convergence
    & same runtime baseline
    & family exists, but no convergence benchmark
    & benchmark pending
    & close \(N=2 \rightarrow 10\) convergence table \\[0.4mm]

    Quadrature-family sensitivity
    & same runtime baseline + quadrature family
    & family exists, but no sensitivity benchmark
    & benchmark pending
    & close Gauss/Lobatto/Radau comparison suite \\
    \bottomrule
  \end{tabularx}
\end{table}

\begin{figure}[htbp]
  \centering
  \begin{tikzpicture}[
      node distance=8mm and 11mm,
      box/.style={draw, rounded corners, align=center, fill=blue!5,
                  minimum width=33mm, minimum height=9mm},
      ready/.style={draw, rounded corners, align=center, fill=green!8,
                    minimum width=33mm, minimum height=9mm},
      pend/.style={draw, rounded corners, align=center, fill=orange!10,
                   minimum width=33mm, minimum height=9mm},
      miss/.style={draw, rounded corners, align=center, fill=red!8,
                   minimum width=33mm, minimum height=9mm},
      arr/.style={-{Latex[length=2.5mm]}, thick}
    ]
    \node[box] (matrix) {Structural matrix\\\(\mathfrak{B}_{\mathrm{phase3}}\)};
    \node[ready, right=of matrix] (runtime) {Runtime-ready claims\\\(C_{\mathrm{rt}}, C_{\sigma_N}, C_{\mathrm{hys}}, C_{M\kappa}\)};
    \node[pend, below=of runtime] (bench) {Benchmark-pending claims\\\(C_N, C_{\mathrm{qsens}}\)};
    \node[pend, below=of matrix] (inst) {Prebenchmark caveat\\base-side station \(\neq\) exact boundary in general};
    \node[box, above=of runtime] (enabler) {Frozen enabler\\\(C_q\)};

    \draw[arr] (matrix) -- (runtime);
    \draw[arr] (matrix) -- (inst);
    \draw[arr] (runtime) -- (bench);
    \draw[arr] (enabler) -- (runtime);
  \end{tikzpicture}
  \caption{The reduced RC column reboot now distinguishes between
  runtime-ready claims and benchmark-pending claims while making explicit that
  the moment-curvature observable now lives at the active base-side station,
  not necessarily at the exact boundary section.  This prevents the structural
  campaign from inheriting physical validation status from a mere smoke run.}
  \label{fig:reduced-rc-column-claim-trace}
\end{figure}

The practical reading of Table~\ref{tab:reduced-rc-column-claim-trace} is
strict.  No row is currently allowed to anchor a physical benchmark by itself.
The catalog deliberately separates three states that are often blurred in
validation work:
\begin{enumerate}
  \item the numerical surface exists and runs;
  \item the corresponding observable exists and is exported cleanly;
  \item the scientific claim supported by that observable has been benchmarked
    against theory, refinement, or experiment.
\end{enumerate}
At the present stage of the reboot, the reduced RC column has reached
State~1 for the full small-strain beam family and State~2 for both the global
hysteresis and the base-side moment-curvature observables.  What remains open
is State~3: benchmark closure.  In particular, the moment-curvature observable
must still be interpreted against the quadrature family because the controlling
station is the nearest active fixed-side section, not a quadrature-independent
ideal boundary cut.  That is precisely why the campaign remains
\emph{pre-benchmark} rather than already validated.

\begin{lstlisting}[language=C++,caption={Canonical reduced-column claim matrix used by the reboot.},label={lst:reduced-rc-column-claim-trace}]
constexpr auto claims =
    fall_n::canonical_reduced_rc_column_validation_claim_table_v;

static_assert(
    fall_n::canonical_reduced_rc_column_claim_runtime_baseline_anchor_count_v == 5);
static_assert(
    fall_n::canonical_reduced_rc_column_claim_instrumentation_pending_count_v == 0);
static_assert(
    fall_n::canonical_reduced_rc_column_claim_physical_benchmark_anchor_count_v == 0);
\end{lstlisting}

\section{Reduced RC column claim-to-benchmark closure matrix}

The claim matrix above is still not the full validation object.  It freezes
what the code is allowed to claim and what evidence exists today, but it does
not yet freeze the \emph{benchmark obligation} that remains attached to each
open claim.  That extra bridge is now kept canonically in
\texttt{src/validation/ReducedRCColumnBenchmarkTraceCatalog.hh}.

For the reduced structural column, the benchmark-trace object can be written
schematically as
\begin{equation}
  \mathfrak{T}_{\mathrm{RC},\,\mathrm{bench}}
  =
  \left\{
    \bigl(C_i,\;B_i,\;\mathcal{O}_i,\;\mathcal{R}_i,\;G_i\bigr)
  \right\}_{i=1}^{6},
\end{equation}
where \(C_i\) is an open reduced-column claim, \(B_i\) the benchmark family
that must be executed, \(\mathcal{O}_i\) the primary observable used for the
comparison, \(\mathcal{R}_i\) the corresponding reference class, and \(G_i\)
the acceptance gate that would close if the benchmark passed.  This follows
the same logic already fixed by the continuum and constitutive chapters:
claims are not validated by existence alone, but by comparing an observable
derived from an energetically consistent discrete object against a declared
reference class \cite{Crisfield1991_NLFEA_Vol1,DeBorst2012_NonlinearFEM}.

\begin{table}[htbp]
  \centering
  \caption{Claim-to-benchmark closure matrix for the reduced RC column reboot.}
  \label{tab:reduced-rc-column-benchmark-trace}
  \scriptsize
  \setlength{\tabcolsep}{3pt}
  \begin{tabularx}{\textwidth}{@{}p{0.19\textwidth}p{0.19\textwidth}p{0.17\textwidth}p{0.17\textwidth}X@{}}
    \toprule
    \textbf{Open claim} & \textbf{Benchmark family} & \textbf{Primary observable}
    & \textbf{Reference class} & \textbf{Acceptance gate closed if passed} \\
    \midrule
    Runtime reduced-column surface
    & runtime-surface matrix
    & base shear vs drift envelope across \(N\) and quadrature families
    & internal runtime-matrix consistency
    & every audited small-strain baseline slice completes with finite response
      and auditable output contracts; this closes the claim itself but is not
      yet a physical benchmark gate. \\[0.4mm]

    Optional axial-compression path
    & axial-load interaction suite
    & base shear vs drift under constant axial preload
    & audited experimental or literature datasets with reported axial-load level
    & axial-load sensitivity becomes eligible to act as part of the phase-3
      structural physical benchmark. \\[0.4mm]

    Hysteresis observable
    & cyclic hysteresis protocol suite
    & base shear vs drift loops
    & audited experimental or literature hysteresis envelopes
    & loop shape, strength degradation, and cycle dissipation become eligible
      to anchor the phase-3 structural physical benchmark. \\[0.4mm]

    Base-side moment-curvature observable
    & base-side section benchmark suite
    & \(\mathcal{O}_{M\kappa}^{\mathrm{base}}\)
    & independent section-level moment-curvature baseline
    & the section observable ceases to be only a runtime contract and becomes
      an audited structural benchmark channel. \\[0.4mm]

    Node-refinement convergence
    & node-refinement suite
    & envelope strength/stiffness/dissipation versus \(N\)
    & internal runtime matrix interpreted as a convergence study
    & refinement no longer remains an unquantified aspiration; it becomes a
      closed precondition for the phase-3 structural benchmark. \\[0.4mm]

    Quadrature-family sensitivity
    & quadrature-sensitivity suite
    & response spread across Gauss/Lobatto/Radau families
    & internal runtime matrix interpreted as a controlled section-placement study
    & quadrature dependence is bounded and made explicit instead of being
      hidden as numerical noise. \\
    \bottomrule
  \end{tabularx}
\end{table}

\begin{figure}[htbp]
  \centering
  \begin{tikzpicture}[
      node distance=8mm and 10mm,
      box/.style={draw, rounded corners, align=center, fill=blue!5,
                  minimum width=33mm, minimum height=9mm},
      bench/.style={draw, rounded corners, align=center, fill=green!8,
                    minimum width=34mm, minimum height=9mm},
      ref/.style={draw, rounded corners, align=center, fill=orange!10,
                  minimum width=35mm, minimum height=9mm},
      gate/.style={draw, rounded corners, align=center, fill=red!8,
                   minimum width=36mm, minimum height=9mm},
      arr/.style={-{Latex[length=2.5mm]}, thick}
    ]
    \node[box] (claim) {Open claim\\\(C_i\)};
    \node[bench, right=of claim] (bench) {Benchmark family\\\(B_i\)};
    \node[ref, right=of bench] (ref) {Reference class\\\(\mathcal{R}_i\)};
    \node[gate, below=of bench] (gate) {Acceptance gate\\\(G_i\)};

    \draw[arr] (claim) -- (bench);
    \draw[arr] (bench) -- (ref);
    \draw[arr] (bench) -- (gate);
    \draw[arr] (ref.south) |- (gate.east);
  \end{tikzpicture}
  \caption{The reduced-column reboot is now traceable one step further:
  every open claim is attached to a concrete benchmark family, a declared
  reference class, and an explicit gate.  This prevents the campaign from
  turning smoke runs or CSV contracts into de facto validation evidence.}
  \label{fig:reduced-rc-column-benchmark-trace}
\end{figure}

Two consequences of Table~\ref{tab:reduced-rc-column-benchmark-trace} are worth
stating explicitly.
\begin{enumerate}
  \item The runtime-surface benchmark is \emph{not} itself a physical benchmark
    gate.  Its role is to prove that the auditable baseline family really
    exists over the declared \(N \times q\) matrix.
  \item The base-side moment-curvature benchmark is intentionally tied first to
    an independent section baseline rather than directly to a full structural
    experiment.  This is scientifically stricter: it isolates the beam-column
    reduction error from the later physical discrepancy between model and test.
\end{enumerate}

\begin{lstlisting}[language=C++,caption={Canonical reduced-column benchmark trace used by the reboot.},label={lst:reduced-rc-column-benchmark-trace}]
constexpr auto trace =
    fall_n::canonical_reduced_rc_column_benchmark_trace_table_v;

static_assert(trace.size() == 6);
static_assert(
    fall_n::canonical_reduced_rc_column_benchmark_external_reference_count_v == 2);
static_assert(
    fall_n::canonical_reduced_rc_column_benchmark_section_baseline_count_v == 1);
static_assert(
    fall_n::canonical_reduced_rc_column_phase3_structural_benchmark_gate_count_v == 5);
\end{lstlisting}

\section{Reduced RC column evidence-closure matrix}

The benchmark-trace layer still answers only \emph{which} benchmark family must
be executed.  A final layer is required before the campaign can be run as a
scientific program instead of a roadmap:
\begin{quote}
  \emph{What evidence artifact already exists, what concrete numerical
  experiment is still missing, and what closure artifact would transform the
  open claim into benchmark evidence?}
\end{quote}
This additional bridge is now frozen canonically in
\texttt{src/validation/ReducedRCColumnEvidenceClosureCatalog.hh}.

For the reduced RC column, the evidence-closure object is represented as
\begin{equation}
  \mathfrak{E}_{\mathrm{RC}}
  =
  \left\{
    \bigl(C_i,\;B_i,\;A_i^{\mathrm{now}},\;X_i^{\mathrm{miss}},\;A_i^{\mathrm{close}}\bigr)
  \right\}_{i=1}^{6},
\end{equation}
where \(C_i\) is the open claim, \(B_i\) the benchmark family already declared
in \(\mathfrak{T}_{\mathrm{RC},\,\mathrm{bench}}\), \(A_i^{\mathrm{now}}\) the
current computational artifact, \(X_i^{\mathrm{miss}}\) the missing numerical
experiment or benchmark dataset, and \(A_i^{\mathrm{close}}\) the closure
artifact that would count as evidence.  This is the strict continuation of the
logic established in Chapters~4--6: an observable is not merely exported, but
must be connected to a reproducible evidence bundle
\cite{Crisfield1991_NLFEA_Vol1,DeBorst2012_NonlinearFEM}.

\begin{table}[htbp]
  \centering
  \caption{Evidence-closure matrix for the reduced RC column reboot.}
  \label{tab:reduced-rc-column-evidence-closure}
  \scriptsize
  \setlength{\tabcolsep}{3pt}
  \begin{tabularx}{\textwidth}{@{}p{0.18\textwidth}p{0.20\textwidth}p{0.20\textwidth}p{0.18\textwidth}X@{}}
    \toprule
    \textbf{Open claim} & \textbf{Current artifact} & \textbf{Missing experiment}
    & \textbf{Closure artifact} & \textbf{Acceptance interpretation} \\
    \midrule
    Runtime reduced-column surface
    & runtime matrix plus audited baseline CSV contracts
    & execute the full \(N \times q\) sweep and summarize finite completion
    & runtime-surface summary table and manifest
    & every declared baseline slice exists and exports the required outputs;
      this closes a runtime completeness obligation, not yet the physical
      benchmark. \\[0.4mm]

    Optional axial-compression path
    & axial-preload branch already running
    & sweep constant axial-load ratios and compare strength/stiffness/dissipation
    & axial-load interaction summary bundle
    & axial-load sensitivity becomes backed by benchmark evidence rather than
      by a single finite smoke run. \\[0.4mm]

    Global hysteresis observable
    & audited hysteresis CSV contract
    & execute the cyclic protocol and compare loop shape, degradation, and
      dissipation against audited reference loops
    & hysteresis comparison bundle
    & the global loop observable becomes a benchmark channel instead of only a
      reproducible export contract. \\[0.4mm]

    Base-side moment-curvature observable
    & \texttt{moment\_curvature\_base.csv}
    & build an independent section-level \(M\)--\(\kappa\) baseline with the
      same section and axial preload
    & section-baseline comparison bundle
    & the structural observable is first validated against a section baseline,
      isolating structural-reduction error before later experiment-to-model
      discrepancy. \\[0.4mm]

    Node-refinement convergence
    & runtime family over \(N=2,\dots,10\)
    & reduce the family to convergence tables for stiffness, strength, and
      dissipation
    & node-refinement convergence bundle
    & refinement becomes a quantified property of the reduced model family
      instead of an assumed qualitative tendency. \\[0.4mm]

    Quadrature-family sensitivity
    & explicit Gauss/Lobatto/Radau model axis
    & reduce the family sweep to spread tables for loops and base-side
      moment--curvature
    & quadrature-sensitivity bundle
    & family dependence is bounded and interpreted explicitly rather than
      hidden under the label of numerical noise. \\
    \bottomrule
  \end{tabularx}
\end{table}

\begin{figure}[htbp]
  \centering
  \begin{tikzpicture}[
      node distance=8mm and 10mm,
      claim/.style={draw, rounded corners, align=center, fill=blue!5,
                    minimum width=29mm, minimum height=9mm},
      art/.style={draw, rounded corners, align=center, fill=green!8,
                  minimum width=33mm, minimum height=9mm},
      miss/.style={draw, rounded corners, align=center, fill=orange!12,
                   minimum width=35mm, minimum height=9mm},
      close/.style={draw, rounded corners, align=center, fill=red!8,
                    minimum width=34mm, minimum height=9mm},
      arr/.style={-{Latex[length=2.5mm]}, thick}
    ]
    \node[claim] (claim) {Open claim\\\(C_i\)};
    \node[art, right=of claim] (now) {Current artifact\\\(A_i^{\mathrm{now}}\)};
    \node[miss, right=of now] (miss) {Missing experiment\\\(X_i^{\mathrm{miss}}\)};
    \node[close, below=of miss] (close) {Closure artifact\\\(A_i^{\mathrm{close}}\)};

    \draw[arr] (claim) -- (now);
    \draw[arr] (now) -- (miss);
    \draw[arr] (miss) -- (close);
  \end{tikzpicture}
  \caption{The reduced-column campaign now carries a fully explicit
  evidence-closure layer.  This is stricter than a benchmark label alone:
  it says what exists today, what still has to be executed, and what artifact
  would count as closure.}
  \label{fig:reduced-rc-column-evidence-closure}
\end{figure}

Two methodological consequences follow directly.
\begin{enumerate}
  \item The current runtime artifact is not yet evidence of validation.  It is
    only the audited computational substrate from which evidence may later be
    produced.
  \item The same benchmark family can require different evidence classes:
    internal runtime summaries, independent section baselines, or external
    datasets.  The reboot therefore forbids collapsing them into one generic
    label such as ``benchmark completed''.
\end{enumerate}

\begin{lstlisting}[language=C++,caption={Canonical reduced-column evidence-closure matrix used by the reboot.},label={lst:reduced-rc-column-evidence-closure}]
constexpr auto evidence =
    fall_n::canonical_reduced_rc_column_evidence_closure_table_v;

static_assert(evidence.size() == 6);
static_assert(
    fall_n::canonical_reduced_rc_column_evidence_external_dataset_count_v == 2);
static_assert(
    fall_n::canonical_reduced_rc_column_evidence_section_baseline_count_v == 1);
static_assert(
    fall_n::canonical_reduced_rc_column_evidence_phase3_gate_count_v == 5);
\end{lstlisting}

\section{Why some attractive extensions remain conditional}

The reboot is intentionally critical with respect to architectural ambition.
Several features are scientifically attractive:
\begin{itemize}
  \item alternative cyclic concrete models or continuum damage-plasticity
    variants \cite{LeeFenves1998_PlasticDamageModelCyclicConcrete,KoBathe2026};
  \item richer reinforced-concrete section and confinement descriptions
    \cite{Mander1988_ConfinedConcrete};
  \item flexibility-based beam-column elements and plastic-hinge integration
    schemes \cite{Spacone1996_FibreBeamColumnModel,Scott2008_SoftwarePatternsBeamColumnModels};
  \item richer Timoshenko beam section placement rules and quadrature families
    \cite{Li2013_FiberSectionTimoshenkoElementOpenSees}.
\end{itemize}

However, scientific precision requires a harder distinction:
\begin{quote}
  \emph{valuable future capability} is not the same as
  \emph{current blocker to the first reference validation campaign}.
\end{quote}
The present plan therefore implements such features only when the audited
campaign shows that the current baseline cannot close its gates honestly.

\section{The first coding artifact of the reboot}

The catalog itself is meant to stay lightweight and outside the hot numerical
path:

\begin{lstlisting}[language=C++,caption={Representative compile-time contract for the rebooted validation campaign.},label={lst:validation-reboot-catalog}]
constexpr auto plan = fall_n::canonical_validation_reboot_workstream_table_v;

static_assert(
    fall_n::canonical_validation_workstream_priority_count_v<
        fall_n::ValidationWorkstreamPriorityKind::mandatory_blocker> == 8);
static_assert(
    fall_n::canonical_validation_reference_structural_column_gate_count_v == 5);
static_assert(
    fall_n::canonical_validation_reference_continuum_column_gate_count_v == 7);
\end{lstlisting}

\section{Immediate execution policy}

The immediate execution policy of this reboot is therefore:
\begin{enumerate}
  \item stabilize the methodological baseline;
  \item validate the reduced-order rectangular RC column first;
  \item build the equivalent continuum column second;
  \item compare both routes quantitatively;
  \item only then reopen walls, frames, subassemblages, FE\textsuperscript{2},
    and building-scale simulations as normative validation targets.
\end{enumerate}

This is narrower than the current ambition of the repository, but also more
scientifically defensible.  That tradeoff is intentional.
